\scanpage{72}
\mspage{- 3 -}
\margin{n° 384}
du support de $\alpha$. Soit $V$ un voisinage de $x$ et soit $A$ un élément de $\Phi$. Soit $\varphi$ une fonction de $\mathcal K_+(G)$, de support contenu dans $V$ et telle que $\alpha(\varphi)\ne 0$. Puisque $\Phi$ converge vers $\alpha$, il existe un élément $B$ de $\Phi$ tel que $\beta(\varphi)\ne 0$ pour $\beta\in B$. Par suite l'ensemble des $\beta\in A$ tels que $V\cap H_\beta\ne\varnothing$ n'est pas vide. Choisissons un tel élément $\beta(V,A)$ et choisissons un point $x(V,A)\in V\cap H_{\beta(V,A)}$. L'application $(V,A)\longrightarrow\beta(V,A)$ (resp. $x(V,A)$) converge vers $\alpha$ (resp. $x$) suivant le filtre des sections du filtre produit du filtre des voisinages de $x$ et du filtre $\Phi$ (ouf ! Vive les directed sets !). D'autre part, on a :
\[
(4)\qquad\beta(V,A)=\beta(V,A)*\varepsilon_{x(V,A)}.
\]
Mais les mesures $\varepsilon_{x(V,A)}$ décrivent un ensemble équicontinu $\mathcal E'(G)$ lorsque $V$ décrit l'ensemble des voisinages de $x$ contenus dans un même voisinage compact. Compte tenu de ce que la topologie vague coïncide sur $\mathfrak M_+(G)$ (resp. sur un ensemble équicontinu de $\mathcal E'(G)$) avec la topologie de la convergence uniforme sur les compacts de $\mathcal K(G)$ (resp. de $\mathcal E(G)$) (Int., ch. III, 2ième édit., § 1, $\ldots$), on déduit alors de (4) et de l'hypocontinuité de la convolution (§ 3, cor. à la prop. 12) que :
\[
(5)\qquad\alpha=\alpha*\varepsilon_x.
\]
Par suite, tout point $y$ de support $H_\alpha$ de $\alpha$ appartient aussi au support de $\alpha*\varepsilon_x$, c'est-à-dire à $H_\alpha x$. Par suite, $yx^{-1}$ appartient à $H_\alpha$ et comme $H_\alpha$ n'est pas vide, c'est un \uline{sous-groupe} de $G$. Enfin, (5) entraîne que la mesure induite par $\alpha$ sur $H_\alpha$ est une mesure de Haar à droite sur $H_\alpha$. Par suite, on a bien $\alpha\in T$.

3. \uline{Bases de} $T$.

Soit $\varphi\in\mathcal K_+(G)$, avec $\varphi(e)>0$. Soit $L_\varphi$ l'hyperplan $\mu(\varphi)=1$ dans $\mathfrak M(G)$. Pour tout $\alpha\in T$, la restriction de $\varphi$ au sous-groupe $H_\alpha$ est un élément non nul de $\mathcal K_+(H_\alpha)$ et on a $\alpha(\varphi)>0$. Par suite, l'hyperplan $L_\varphi$ rencontre chaque géné-
\margin{71}
